\chapter{Theoretical Background}

\section{Introduction}
This chapter presents the theoretical background and key technologies relevant to the development of the proposed system. The project combines concepts from computer vision, embedded artificial intelligence, robotics, and software engineering. In particular, it focuses on the challenges associated with deploying machine learning models on resource-constrained devices and integrating them into a cohesive robotic system. The purpose of this chapter is to provide the necessary foundation for understanding the design and implementation decisions described in later chapters.

\section{Computer Vision for Embedded Object Detection}

\subsection{Overview of Computer Vision in Robotics}

Computer vision is a subfield of artificial intelligence that enables machines to extract meaningful information from images and video streams. In robotic systems, it serves as the primary perception mechanism, allowing autonomous or semi-autonomous systems to interpret and react to their environment. The general objective is to transform raw pixel data into structured information that can support decision-making processes.

\subsection{Object Detection in Visual Systems}

Object detection is the computer-vision task of simultaneously
\emph{classifying} and \emph{localising} one or more objects within an
image \cite{zou2023object}. Its output is typically a set of bounding
boxes, each paired with a class label and a confidence score.

Modern detectors are broadly divided into two families:
 
\begin{itemize}
    \item \textbf{Two-stage detectors} (e.g., Faster R-CNN) first
          generate region proposals and then classify each
          proposal~\cite{ren2015faster}. They tend to be highly
          accurate but computationally expensive.
 
    \item \textbf{Single-stage detectors} (e.g., YOLO, SSD) predict
          boxes and classes in a single forward pass~\cite{redmon2016yolo},
          making them well-suited for real-time and embedded
          applications.
\end{itemize}
 
Performance is typically measured with \emph{mean Average Precision}
(mAP), which summarises precision–recall trade-offs across all classes
and Intersection-over-Union (IoU) thresholds~\cite{everingham2010pascal}.

 
% ------------------------------------------------------------
\section{Convolutional Neural Networks}
% ------------------------------------------------------------
 
A Convolutional Neural Network (CNN) is a class of deep neural network
designed specifically for grid-structured data such as
images~\cite{lecun1998gradient}. Three operations form its core:
 
\begin{enumerate}
    \item \textbf{Convolution} — a learnable filter slides over the
          input to extract local spatial features (edges, textures,
          shapes).
    \item \textbf{Pooling} — spatially downsamples feature maps,
          reducing computation and providing translation invariance.
    \item \textbf{Fully connected layers} — aggregate extracted
          features for final classification or regression.
\end{enumerate}
 
CNNs are central to object detection because their hierarchical
feature extraction naturally captures low-level details (edges) in
early layers and high-level semantics (object parts) in deeper
layers~\cite{zeiler2014visualizing}. Backbone networks such as
ResNet~\cite{he2016deep} and MobileNet~\cite{howard2017mobilenets}
are commonly used as feature extractors inside detection pipelines,
with MobileNet-style architectures preferred on embedded hardware due
to their use of depthwise separable convolutions that drastically
reduce parameter count and FLOPs.
 
% ------------------------------------------------------------
\section{Deep Learning}
% ------------------------------------------------------------
 
Deep learning refers to machine learning with neural networks
containing many stacked layers, enabling automatic discovery of
hierarchical representations directly from raw
data~\cite{lecun2015deep}. Training is performed end-to-end via
\emph{backpropagation} and \emph{stochastic gradient descent} (SGD),
minimising a task-specific loss function over large labelled
datasets~\cite{goodfellow2016deep}.
 
Key factors that made deep learning dominant in computer vision are:
 
\begin{itemize}
    \item Large annotated datasets (ImageNet, COCO).
    \item GPU-accelerated parallel computation.
    \item Architectural innovations: batch normalisation, residual
          connections, attention mechanisms.
\end{itemize}
 
More recently, \emph{transformer}-based architectures, originally
developed for natural language processing~\cite{vaswani2017attention},
have been successfully adapted to vision tasks through the Vision
Transformer (ViT)~\cite{dosovitskiy2021image} and its descendants,
achieving state-of-the-art results by modelling long-range spatial
dependencies via self-attention — a capability CNNs approximate only
indirectly through stacking layers.

\subsection{RF-DETR Nano Model Overview}

RF-DETR (Roboflow Detection Transformer) is a real-time
transformer-based object detection architecture developed by
Roboflow~\cite{rfdetr2024roboflow}. It is built on a
\textbf{DINOv2}~\cite{oquab2023dinov2} vision transformer backbone,
a self-supervised ViT pre-trained on large-scale image data that
produces rich, generalizable feature representations without
task-specific supervision.
 
The architecture follows the DETR paradigm~\cite{carion2020end}:
a transformer encoder–decoder processes backbone features and
predicts a fixed set of object queries, each decoded into a bounding
box and class label — eliminating hand-crafted components such as
anchor generation and non-maximum suppression (NMS).
 
The \textbf{Nano} variant is a size-optimised configuration of
RF-DETR designed for deployment in resource-constrained environments.
It reduces the number of encoder layers, embedding dimensions, and
decoder queries relative to larger variants, achieving a favourable
accuracy–latency trade-off suitable for embedded hardware.
 
Key characteristics of RF-DETR Nano relevant to this work:
 
\begin{itemize}
    \item \textbf{Backbone:} DINOv2 ViT (small configuration),
          providing strong pre-trained features.
    \item \textbf{Decoder:} Transformer decoder with learnable object
          queries; end-to-end trained with a Hungarian matching
          loss~\cite{carion2020end}.
    \item \textbf{Real-time inference:} Optimised for low-latency
          detection without sacrificing significant accuracy compared
          to larger models.
    \item \textbf{No NMS:} The set-prediction formulation removes the
          post-processing bottleneck common in CNN-based detectors.
\end{itemize}
 
Table~\ref{tab:rfdetr_compare} situates RF-DETR Nano among comparable
lightweight detectors.
 
\begin{table}[h]
\centering
\caption{Comparison of lightweight object detection models.}
\label{tab:rfdetr_compare}
\begin{tabular}{lccc}
\hline
\textbf{Model} & \textbf{Params (M)} & \textbf{mAP (COCO)} & \textbf{Architecture} \\
\hline
YOLOv8n        & 3.2  & 37.3 & CNN            \\
RT-DETR-R18    & 20.0 & 46.5 & Transformer    \\
RF-DETR Nano   & 30.5   & 48.4   & ViT + DETR     \\
\hline
\end{tabular}
\end{table}

\section{Model Deployment and Format Conversion for Edge Devices}
% ------------------------------------------------------------
 
Deploying deep learning models on embedded systems requires more than
simply porting trained weights to a new device. Standard training
frameworks such as PyTorch or TensorFlow are designed for
high-performance workstations with abundant memory and floating-point
units; they carry significant runtime overhead that is impractical on
resource-constrained hardware~\cite{warden2019tinyml}. Bridging this
gap involves two complementary strategies: \emph{format conversion},
which re-expresses the model in a hardware-friendly representation,
and \emph{quantization}, which reduces numerical precision to lower
memory and compute requirements. Both are explored in this section.
 
% ............................................................
\subsection{ONNX (Open Neural Network Exchange)}
% ............................................................
 
ONNX is an open, vendor-neutral format for representing neural network
models~\cite{bai2019onnx}. It defines a common computational graph
whose nodes map to standard mathematical operations (convolutions,
matrix multiplications, activations), allowing a model trained in
one framework — say PyTorch — to be exported once and executed by
any ONNX-compatible runtime without rewriting the model code.
 
Key properties that make ONNX attractive for deployment pipelines:
 
\begin{itemize}
    \item \textbf{Framework interoperability.} Models from PyTorch,
          TensorFlow, and others can all be serialised to the same
          \texttt{.onnx} file and consumed by runtimes such as
          ONNX Runtime, OpenVINO, or TensorRT.
 
    \item \textbf{Graph optimisation.} ONNX Runtime applies
          graph-level rewrites (operator fusion, constant folding)
          that reduce inference latency independently of the original
          framework~\cite{onnxruntime2023}.
 
    \item \textbf{Hardware portability.} The same \texttt{.onnx} file
          can target CPUs, GPUs, and dedicated accelerators through
          pluggable \emph{execution providers}, requiring no changes
          to the model graph itself.
\end{itemize}
 
In the context of embedded deployment, ONNX serves as an intermediate
step: the model is first exported to ONNX to decouple it from the
training framework, and then either executed directly via ONNX Runtime
or converted further into a hardware-specific format.
 
% ............................................................
\subsection{TensorFlow Lite}
% ............................................................
 
TensorFlow Lite (TFLite) is a dedicated inference framework developed
by Google for mobile and embedded devices~\cite{david2021tensorflow}.
Unlike ONNX, which is primarily a \emph{model format}, TFLite is a
complete \emph{inference stack} — it defines its own FlatBuffer-based
model format (\texttt{.tflite}), a lightweight interpreter, and a
set of hardware-specific delegates.
 
\begin{itemize}
    \item \textbf{FlatBuffer serialisation.} The \texttt{.tflite}
          format is memory-mapped at runtime, allowing the interpreter
          to begin inference without loading the entire model into
          RAM — a critical advantage on devices with limited
          memory~\cite{david2021tensorflow}.
 
    \item \textbf{Delegate system.} Delegates offload portions of the
          graph to accelerators: the GPU delegate targets OpenGL ES /
          Metal, the NNAPI delegate targets Android neural network
          hardware, and the XNNPACK delegate accelerates
          floating-point operations on ARM CPUs via SIMD
          intrinsics~\cite{xnnpack2020}.
 
    \item \textbf{Built-in quantization support.} TFLite provides
          first-class tooling for post-training quantization and
          quantization-aware training (discussed in
          Section~\ref{sec:quantization}), making it the most common
          end-to-end pipeline for deploying quantised models on
          edge hardware.
\end{itemize}
 
% ............................................................
\subsection{Trade-offs in Format Conversion}
% ............................................................
 
Converting a model from its training format to a deployment format is
rarely lossless. Practitioners must navigate several trade-offs:
 
\begin{itemize}
    \item \textbf{Operator coverage.} Not every layer or activation
          function is supported by every runtime. Custom or exotic
          operators may be unsupported, requiring architectural
          changes before conversion succeeds.
 
    \item \textbf{Numerical fidelity.} Graph optimisations such as
          operator fusion change the order of floating-point
          operations, which can introduce small numerical differences
          compared to the original framework's output.
 
    \item \textbf{Latency vs.\ accuracy.} Aggressively optimised
          formats (especially when combined with quantization) reduce
          inference time at the cost of a measurable drop in
          prediction accuracy. Profiling on the target hardware is
          essential to find an acceptable operating point.
 
    \item \textbf{Toolchain complexity.} Each conversion step
          (framework $\to$ ONNX $\to$ TFLite, for example) introduces
          a potential failure point and requires version-compatible
          tools.
\end{itemize}
 
In this project, both ONNX and TensorFlow Lite formats were evaluated
as deployment candidates for the RF-DETR Nano model on the Raspberry Pi
target platform.
 
% ............................................................
\subsection{Model Quantization}
\label{sec:quantization}
% ............................................................
 
Quantization is the process of representing model weights and/or
activations with lower-precision numeric types than the 32-bit
floating-point (FP32) format used during training~\cite{jacob2018quantization}.
It is one of the most effective techniques for reducing model size and
accelerating inference on hardware that lacks native FP32 or FP16
units — a common situation in microcontrollers and low-power
SoCs~\cite{krishnamoorthi2018quantizing}.
 
\subsubsection{Floating-Point Precision Levels}
 
Before discussing quantization schemes it is helpful to understand the
numeric formats involved:
 
\begin{itemize}
    \item \textbf{FP32 (32-bit float).} The IEEE 754 single-precision
          standard: 1 sign bit, 8 exponent bits, 23 mantissa bits.
          Used by default during training. High dynamic range and
          precision but expensive in memory and compute.
 
    \item \textbf{FP16 (16-bit float).} Half-precision: 1 sign,
          5 exponent, 10 mantissa bits. Halves memory and bandwidth
          requirements; natively accelerated on modern GPUs and some
          edge NPUs. Introduces a reduced dynamic range that can cause
          gradient underflow during training but is generally
          acceptable for inference~\cite{micikevicius2018mixed}.
 
    \item \textbf{BF16 (bfloat16).} An alternative 16-bit format with
          the same 8-bit exponent as FP32 but only 7 mantissa bits.
          Retains FP32's dynamic range, making it numerically safer
          for training on TPUs and certain AI accelerators.
 
    \item \textbf{INT8 (8-bit integer).} No fractional bits; values
          mapped to the integer range $[-128, 127]$ (signed) or
          $[0, 255]$ (unsigned) via a learned \emph{scale} and
          \emph{zero-point}. Reduces model size by $4\times$ versus
          FP32 and enables fast integer multiply-accumulate (MAC)
          operations available on virtually all embedded
          processors~\cite{jacob2018quantization}.
 
    \item \textbf{INT4 / binary (1-bit).} Ultra-low precision formats
          used in research contexts. They offer extreme compression
          but typically incur significant accuracy loss and have
          limited hardware support~\cite{hubara2017quantized}.
\end{itemize}
 
\subsubsection{Types of Quantization}
 
There are two principal strategies for quantizing a trained model:
 
\paragraph{Post-Training Quantization (PTQ).}
PTQ converts a pre-trained FP32 model to lower precision without
retraining~\cite{nagel2020up}. A small \emph{calibration dataset}
(typically a few hundred representative images) is passed through
the model to record the statistical range of each layer's activations.
These statistics are used to compute the scale and zero-point
parameters needed to map floating-point values to integers. PTQ
is fast and convenient — no labelled training data are required —
but can degrade accuracy noticeably on models with wide activation
distributions.
 
Two sub-variants are common in practice:
 
\begin{itemize}
    \item \textbf{Dynamic quantization.} Only the weights are
          quantized ahead of time; activations are quantized
          on-the-fly during inference. Straightforward to apply
          but offers less speedup than full static quantization.
 
    \item \textbf{Static quantization.} Both weights and activations
          are quantized using calibration statistics gathered
          offline. Faster at runtime but requires a representative
          calibration set and careful range selection.
\end{itemize}
 
\paragraph{Quantization-Aware Training (QAT).}
QAT simulates quantization during the forward pass of training by
inserting \emph{fake quantization} nodes — operations that round
values to the nearest representable integer and then immediately
dequantize them back to float~\cite{jacob2018quantization}. Because
gradients flow through these nodes, the model learns to compensate
for quantization error. QAT consistently produces higher accuracy than
PTQ, particularly at INT8, but requires access to training data and
a longer training pipeline.
 
Table~\ref{tab:quant_compare} summarises the trade-offs between the
two approaches.
 
\begin{table}[h]
\centering
\caption{Comparison of quantization strategies.}
\label{tab:quant_compare}
\begin{tabular}{lccc}
\hline
\textbf{Strategy} & \textbf{Training Required} & \textbf{Accuracy} & \textbf{Effort} \\
\hline
Post-Training (Dynamic)  & No  & Moderate & Low    \\
Post-Training (Static)   & No  & Good     & Low    \\
Quantization-Aware (QAT) & Yes & Best     & High   \\
\hline
\end{tabular}
\end{table}
 
% ------------------------------------------------------------
\section{Embedded Systems for AI and Robotics}
% ------------------------------------------------------------
 
An \emph{embedded system} is a dedicated computing platform integrated
into a larger device to perform a specific, well-defined
function~\cite{heath2002embedded}. Unlike general-purpose computers,
embedded systems are optimised for a particular workload and operate
under tight constraints of processing power, memory, storage, and
energy consumption.
 
% ............................................................
\subsection{Edge Computing}
% ............................................................
 
\emph{Edge computing} refers to the execution of computation at or
near the source of data — on the device itself — rather than
offloading to a remote cloud server~\cite{shi2016edge}. For
robotics and real-time vision applications this distinction is
critical:
 
\begin{itemize}
    \item \textbf{Latency.} A round trip to a cloud API can add
          tens to hundreds of milliseconds of delay, which is
          unacceptable for control loops that must respond within
          milliseconds~\cite{shi2016edge}.
 
    \item \textbf{Reliability.} On-device inference continues to
          function when network connectivity is unavailable or
          intermittent — common in agricultural or field robotics
          scenarios.
 
    \item \textbf{Privacy and bandwidth.} Processing sensor data
          locally avoids transmitting raw images or video streams,
          reducing both bandwidth cost and data-privacy risk.
\end{itemize}
 
% ............................................................
\subsection{Constraints of Embedded AI}
% ............................................................
 
Running inference on embedded hardware is fundamentally different from
running it on a workstation. Three constraints dominate system
design~\cite{warden2019tinyml}:
 
\begin{itemize}
    \item \textbf{Compute.} Embedded CPUs operate at lower clock
          frequencies and lack the massively parallel floating-point
          units of desktop GPUs. Operations that are trivially fast on
          a workstation may become the bottleneck on a single-board
          computer.
 
    \item \textbf{Memory.} RAM is typically measured in hundreds of
          megabytes rather than gigabytes. Large model activations or
          feature maps may exceed available memory, requiring
          architectural choices — smaller input resolution,
          fewer layers, or quantization — to bring the memory
          footprint within bounds.
 
    \item \textbf{Power.} Battery-operated platforms impose strict
          power budgets. Higher computational load increases current
          draw, reducing operating time and generating heat that
          must be managed passively.
\end{itemize}
 
These constraints motivate every optimisation technique discussed in
Section~\ref{sec:quantization}: reducing precision from FP32 to INT8
lowers memory by $4\times$, reduces bandwidth, and replaces
power-hungry floating-point MAC operations with cheaper integer ones.
 
% ............................................................
\subsection{The Raspberry Pi as an Edge Computing Platform}
% ............................................................
 
The Raspberry Pi is a single-board computer (SBC) developed by the
Raspberry Pi Foundation~\cite{upton2016raspberry}. It integrates an
ARM-based SoC, RAM, USB, GPIO, and camera interfaces onto a small,
low-cost PCB, making it one of the most widely adopted platforms for
prototyping embedded AI systems.
 
In this project, the Raspberry Pi serves as the primary inference
platform. It offers a Linux environment that supports standard Python
tooling, ONNX Runtime, and TensorFlow Lite, while remaining
representative of the class of devices targeted in real-world
edge deployments. Its hardware limitations — a multi-core ARM CPU
with no dedicated neural-network accelerator — make it an appropriate
testbed for evaluating the latency impact of format conversion and
quantization on RF-DETR Nano.


